% problem-000062 4 乙酰丙酮光度测醛
\section*{4. 乙酰丙酮光度测醛}
\subsection*{4-1 $\mathrm { N a _ { 2 } S _ { 2 } O _ { 3 } }$ 溶液的标定。}

移取 25.00 cm3 浓度为 $8 . 3 3 3 \times 1 0 ^ { - 3 } \mathrm { m o l d m ^ { 3 } }$ 的 $\mathrm { K _ { 2 } C r _ { 2 } O _ { 7 } }$ _{7}标准溶液于 $2 5 0 \mathrm { c m } ^ { 3 }$ 碘量瓶中，加⼊ KI，⽤硫酸酸化，⾄反应完全，⽣成I_{3}-，⽤ $\mathrm { N a _ { 2 } S _ { 2 } O _ { 3 } }$ 溶液滴定，消耗 $2 4 . 2 2 \mathrm { c m } ^ { 3 } \mathrm { c m } ^ { 3 } \mathrm { c m } $ 。计算 $\mathrm { N a _ { 2 } S _ { 2 } O _ { 3 } }$ 溶液的浓度。

\subsection*{4-2 甲醛标准储备液的标定。}

移取 $2 0 . 0 0 \mathrm { c m } ^ { 3 }$ 甲醛标准储备液于 $2 5 0 \mathrm { c m } ^ { 3 }$ 碘量瓶中，加⼊ $5 0 . 0 0 \mathrm { c m } ^ { 3 }$ 浓度为 $0 . 0 2 5 \mathrm { m o l d m ^ { - 3 } }$ 碘溶液和适量NaOH溶液使反应完全，然后加⼊适量稀硫酸酸化⾄弱酸性。⽤上述 $\mathrm { N a _ { 2 } S _ { 2 } O _ { 3 } }$ 标准溶液滴定⾄终点，消耗$2 2 . 0 2 \mathrm { c m } ^ { 3 } \mathrm { c m }$ 。⽤相同步骤做空⽩实验，消耗 $4 8 . 1 0 \mathrm { c m } ^ { 3 } \mathrm { c m }$ 。写出碱性条件下甲醛与碘溶液反应⽅程式，计算甲醛标准储备液的浓度（HCHO，mg cm-3）

\subsection*{4-3 废⽔中甲醛含量的测定。}

将甲醛标准储备液⽤去离⼦⽔逐级稀释100倍得甲醛标准使⽤液。向数⽀ $2 5 \mathrm { c m ^ { 3 } }$ 的⽐⾊管中分别加⼊不同体积的甲醛标准使⽤液，加⽔⾄标线，分别加⼊ $2 . 5 0 \mathrm { c m } ^ { 3 }$ ⼄酰丙酮显⾊剂溶液，混匀，于 $4 5 { \sim } 6 0 ~ ^ { \circ } \mathrm { C }$ ⽔浴中加热30 min，取出冷却。于414 nm波⻓处，以去离⼦⽔作参⽐，⽤1 cm⽐⾊⽫测定吸光度�。

另移取废⽔25.00 cm3的⽐⾊管中，按上述步骤进⾏显⾊和吸光度测量，所得数据如下：

\textit{（原卷此处含表格/仪器试剂表，详见 PDF 或 MD 源文件。）}

计算出标准曲线的回归⽅程，并计算废⽔中甲醛的含量(HCHO, mg dm−3)。

提⽰：标准曲线的回归⽅程中�为直线的截距，�为直线的斜率，其中�与�可按下式求解。

$
b = \frac {\sum_ {i = 1} ^ {n} (x _ {i} - \bar {x}) (y _ {i} - \bar {y})}{\sum_ {i = 1} ^ {n} (x _ {i} - \bar {x}) ^ {2}} a = \bar {y} - b \bar {x}
$
